\section{Probabilistic Circuits}

\begin{frame}{Tensor xác suất và động lực}
\[
    \mathcal{P}\in
    \mathbb{R}^{|\mathcal{D}_1|\times\cdots\times|\mathcal{D}_N|},
    \qquad
    |\mathcal{P}|=\prod_{n=1}^{N}|\mathcal{D}_n|
\]
\begin{itemize}
    \item Mỗi entry là một xác suất $p(x_1,\ldots,x_N)$.
    \item Tensor xác suất đầy đủ tăng theo hàm mũ theo số biến.
    \item Tensor factorization biểu diễn gọn; circuit tổ chức phép tính để truy vấn hiệu quả.
\end{itemize}
\[
    \text{Tensor Factorization}
    \Longleftrightarrow
    \text{Circuit}
    \Longrightarrow
    \text{Tractable Inference}
\]
\end{frame}

\begin{frame}{Circuit là gì?}
\begin{columns}[T,totalwidth=\textwidth]
    \begin{column}{0.46\textwidth}
        \centering
        \begin{tikzpicture}[
            node/.style={circle, draw, thick, minimum size=0.62cm, font=\small},
            input/.style={rectangle, draw, rounded corners=3pt, fill=teal!12, minimum width=1.0cm, minimum height=0.45cm, font=\scriptsize},
            sum/.style={node, fill=blue!12},
            prod/.style={node, fill=orange!14},
            arrow/.style={-{Stealth[length=2.2mm]}, thick}
        ]
            \node[sum] (s) at (0,0) {$+$};
            \node[prod] (p1) at (-1.2,-1.0) {$\times$};
            \node[prod] (p2) at (1.2,-1.0) {$\times$};
            \node[input] (x1) at (-2.0,-2.0) {$f_1(x_1)$};
            \node[input] (x2) at (-0.6,-2.0) {$f_2(x_2)$};
            \node[input] (x3) at (0.6,-2.0) {$g_1(x_1)$};
            \node[input] (x4) at (2.0,-2.0) {$g_2(x_2)$};
            \draw[arrow] (x1) -- (p1);
            \draw[arrow] (x2) -- (p1);
            \draw[arrow] (x3) -- (p2);
            \draw[arrow] (x4) -- (p2);
            \draw[arrow] (p1) -- (s);
            \draw[arrow] (p2) -- (s);
        \end{tikzpicture}
    \end{column}
    \begin{column}{0.50\textwidth}
        \begin{itemize}
            \item Circuit là một DAG tham số hóa, tính hàm bằng forward pass theo thứ tự topo.
            \item Mỗi node có \textbf{scope}: tập biến mà node phụ thuộc.
            \item Kích thước $|c|$ thường tính bằng số node/cạnh/tham số.
        \end{itemize}
        \[
            c(x)=\sum_{j\in\mathrm{in}(n)}w_{n,j}c_j(x)
            \quad\text{hoặc}\quad
            c(x)=\prod_{j\in\mathrm{in}(n)}c_j(x)
        \]
    \end{column}
\end{columns}
\end{frame}

\begin{frame}{Circuit gồm những unit nào?}
\small
\begin{columns}[T,totalwidth=\textwidth]
    \begin{column}{0.32\textwidth}
        \textbf{Input unit}
        \[
            c_n(x_{S_n})=f_n(x_{S_n})
        \]
        \begin{itemize}
            \item Đọc biến hoặc hàm cục bộ.
            \item Scope nhỏ.
        \end{itemize}
    \end{column}
    \begin{column}{0.32\textwidth}
        \textbf{Product unit}
        \[
            c_n=\prod_{j\in\mathrm{in}(n)}c_j
        \]
        \begin{itemize}
            \item Ghép các scope con.
            \item Tractable khi decomposable.
        \end{itemize}
    \end{column}
    \begin{column}{0.32\textwidth}
        \textbf{Sum unit}
        \[
            c_n=\sum_{j\in\mathrm{in}(n)}w_{n,j}c_j
        \]
        \begin{itemize}
            \item Mixture/tổng có trọng số.
            \item Tractable khi smooth.
        \end{itemize}
    \end{column}
\end{columns}
\[
    \mathrm{scope}(n)=\bigcup_{j\in\mathrm{in}(n)}\mathrm{scope}(j)
\]
\end{frame}
\begin{frame}{CP decomposition như một circuit sum-product}
\begin{columns}[T,totalwidth=\textwidth]
    \begin{column}{0.60\textwidth}
        \centering
        \begin{tikzpicture}[
            sum/.style={circle, draw, fill=blue!10, minimum size=0.65cm, thick},
            prod/.style={circle, draw, fill=orange!10, minimum size=0.65cm, thick},
            leaf/.style={rectangle, draw, rounded corners=3pt, fill=teal!10, minimum height=0.45cm, minimum width=1.05cm, thick, font=\scriptsize},
            arrow/.style={-{Stealth[length=2.2mm]}, thick}
        ]
            \node[sum] (out) at (0,0) {$+$};
            \node[prod] (p1) at (-1.8,-1.1) {$\times$};
            \node[prod] (p2) at (1.8,-1.1) {$\times$};
            \foreach \x/\name/\lab in {-2.8/l11/$V^{(1)}_{x_1,1}$,-1.8/l12/$V^{(2)}_{x_2,1}$,-0.8/l13/$V^{(3)}_{x_3,1}$} {
                \node[leaf] (\name) at (\x,-2.4) {\lab};
                \draw[arrow] (\name) -- (p1);
            }
            \foreach \x/\name/\lab in {0.8/l21/$V^{(1)}_{x_1,2}$,1.8/l22/$V^{(2)}_{x_2,2}$,2.8/l23/$V^{(3)}_{x_3,2}$} {
                \node[leaf] (\name) at (\x,-2.4) {\lab};
                \draw[arrow] (\name) -- (p2);
            }
            \draw[arrow] (p1) -- (out);
            \draw[arrow] (p2) -- (out);
        \end{tikzpicture}
    \end{column}
    \begin{column}{0.36\textwidth}
        \[
            \mathcal{T}_{x_1,x_2,x_3}
            =
            \sum_{r=1}^{R}
            V^{(1)}_{x_1,r}
            V^{(2)}_{x_2,r}
            V^{(3)}_{x_3,r}
        \]
        \begin{itemize}
            \item Product nodes: nhân theo mode.
            \item Sum node: cộng theo rank.
        \end{itemize}
    \end{column}
\end{columns}
\end{frame}

\begin{frame}{Tucker circuit và chi phí core}
\[
    \mathcal{T}_{x_1,\ldots,x_N}
    \approx
    \sum_{r_1,\ldots,r_N}
    \mathcal{W}_{r_1,\ldots,r_N}
    \prod_{n=1}^{N}V^{(n)}_{x_n,r_n}
\]
\begin{itemize}
    \item Tucker circuit dùng core tensor để trộn các latent dimensions.
    \item Kích thước circuit/core có thể tăng như:
    \[
        O\!\left(N\prod_{n=1}^{N}R_n\right)
    \]
    \item Khi core superdiagonal, Tucker trở về CP.
\end{itemize}
\end{frame}

\begin{frame}{Structural properties}
\small
\renewcommand{\arraystretch}{1.25}
\begin{tabular}{p{0.24\textwidth}p{0.38\textwidth}p{0.28\textwidth}}
\toprule
\textbf{Property} & \textbf{Điều kiện} & \textbf{Cho phép} \\
\midrule
Smoothness & Các nhánh vào sum node có cùng scope & Marginalization đúng \\
Decomposability & Các nhánh vào product node có scope rời nhau & Tách tích phân theo nhóm biến \\
Compatibility & Hai circuits chia scope cùng cấu trúc & Tính expectation giữa $p$ và $f$ \\
Determinism & Supports của nhánh sum rời nhau & MAP inference hiệu quả hơn \\
\bottomrule
\end{tabular}
\[
    \text{smooth + decomposable}
    \Rightarrow
    \text{marginals trong }O(|c|)
\]
\end{frame}

\begin{frame}{Truy vấn suy luận: marginal và conditional}
\begin{columns}[T,totalwidth=\textwidth]
    \begin{column}{0.48\textwidth}
        \textbf{Marginal}
        \[
            p(\mathbf{x}_O)
            =
            \sum_{\mathbf{x}_M\in\mathcal{D}_M}
            p(\mathbf{x}_O,\mathbf{x}_M)
        \]
        \begin{itemize}
            \item $O$: biến quan sát, $M$: biến bị sum-out.
            \item Smooth + decomposable $\Rightarrow O(|c|)$.
        \end{itemize}
    \end{column}
    \begin{column}{0.48\textwidth}
        \textbf{Conditional}
        \[
            p(\mathbf{x}_M\mid\mathbf{x}_O)
            =
            \frac{p(\mathbf{x}_O,\mathbf{x}_M)}
            {\sum_{\mathbf{x}'_M}p(\mathbf{x}_O,\mathbf{x}'_M)}
        \]
        \begin{itemize}
            \item Tử số: evaluate circuit tại assignment.
            \item Mẫu số: marginalize biến missing.
        \end{itemize}
    \end{column}
\end{columns}
\end{frame}

\begin{frame}{Truy vấn suy luận dưới góc nhìn tensor factorization}
\[
    c(x_1,\ldots,x_N)
    =
    \sum_{r=1}^{R}\prod_{n=1}^{N}V^{(n)}_{x_n,r}
\]
\[
    c(\mathbf{x}_O)
    =
    \sum_{r=1}^{R}
    \left(\prod_{n\in O}V^{(n)}_{x_n,r}\right)
    \left(\prod_{m\in M}\sum_{x_m\in\mathcal{D}_m}V^{(m)}_{x_m,r}\right)
\]
\begin{itemize}
    \item Marginalize một mode $m$ bằng cách thay factor $V^{(m)}_{x_m,r}$ bởi tổng theo hàng của mode đó.
    \item Vì product node decomposable, phép cộng trên nhiều assignment được đẩy xuống từng factor matrix.
\end{itemize}
\end{frame}

\begin{frame}{Truy vấn suy luận: expectation}
\[
    \mathbb{E}_{\mathbf{x}\sim p}[f(\mathbf{x})]
    =
    \sum_{\mathbf{x}}p(\mathbf{x})f(\mathbf{x})
    =
    \langle\mathcal{P},\mathcal{F}\rangle
\]
\begin{itemize}
    \item Nếu $p$ và $f$ là hai circuits compatible, có thể tính bằng product circuit rồi marginalize.
    \item Dưới góc nhìn TF: contract tensor xác suất $\mathcal{P}$ với tensor hàm $\mathcal{F}$.
    \item Ứng dụng trong report: missing values, fairness query, adversarial robustness.
\end{itemize}
\[
    \mathbb{E}[f(\mathbf{x})\mid A=0]-\mathbb{E}[f(\mathbf{x})\mid A=1]
\]
\end{frame}

\begin{frame}{Truy vấn suy luận: MAP và sampling}
\begin{columns}[T,totalwidth=\textwidth]
    \begin{column}{0.48\textwidth}
        \textbf{MAP / MPE}
        \[
            \mathbf{x}^{\star}
            =
            \arg\max_{\mathbf{x}}p(\mathbf{x}\mid\mathbf{e})
        \]
        \begin{itemize}
            \item Cần thêm determinism để tối ưu tractable hơn.
            \item Sum node chọn nhánh tốt nhất thay vì cộng tất cả nhánh.
        \end{itemize}
    \end{column}
    \begin{column}{0.48\textwidth}
        \textbf{Sampling top-down}
        \begin{itemize}
            \item Sum node: lấy mẫu nhánh categorical theo trọng số.
            \item Product node: đi qua tất cả nhánh con vì scope rời nhau.
            \item Leaf/input: lấy mẫu biến quan sát cục bộ.
        \end{itemize}
    \end{column}
\end{columns}
\end{frame}
\begin{frame}{Non-negative TF như latent variable model}
\[
    p(x)
    =
    \sum_{r=1}^{R}p(r)\prod_{n=1}^{N}p_n(x_n\mid r)
\]
\begin{itemize}
    \item Chỉ số rank $r$ có thể đọc như latent component.
    \item Sum node tương ứng với mixture variable.
    \item Product node sinh độc lập trên các scope rời nhau.
    \item Sampling thực hiện top-down: chọn nhánh ở sum, đi qua tất cả nhánh ở product.
\end{itemize}
\end{frame}

\begin{frame}{Demo CPJ Circuit}
\centering
\renewcommand{\arraystretch}{1.25}
\begin{tabular}{p{0.28\textwidth}p{0.30\textwidth}p{0.30\textwidth}}
\toprule
\textbf{Demo} & \textbf{Thiết lập} & \textbf{Kết quả chính} \\
\midrule
Toy 4-bit & $2^4=16$ assignments & KL(target $\Vert$ circuit) $=0.0054$; tổng xác suất $=1.000000$ \\
Binarized MNIST & $28\times 28=784$ biến nhị phân & Best val bpd $=0.3675$; test bpd $=0.3704$ \\
\bottomrule
\end{tabular}
\[
    \mathrm{bpd}=-\frac{\log p(x)}{784\log 2}
\]
\end{frame}

\begin{frame}{Hạn chế và câu hỏi mở của PC}
\begin{itemize}
    \item Tractable inference không đảm bảo biểu diễn luôn nhỏ.
    \item Tucker/core có thể tăng theo $R^N$ nếu số mode lớn.
    \item Hiệu quả phụ thuộc region graph, width/rank và cấu trúc phụ thuộc thật trong dữ liệu.
    \item Câu hỏi mở: optimization, identifiability, approximate inference, scalability.
\end{itemize}
\end{frame}
